\chapter{Актуальные проблемы управления интернет сетей и методы их решения}\label{ch:ch1}

\section{Общие проблемы анализа трафика и сетей}\label{sec:ch1/sec1}

Задачи анализа состояния сети и сетевого трафика, выявление различных свойств его содержимого играют важную роль в современных сетях. Эти задачи являются важным аспектом обеспечения необходимого уровня QoE (Quality of Experience, качество восприятия пользователем) и QoS (Quality of Service, качество обслуживания) за счет предоставления различным классам трафика различных приоритетов в обслуживании. Также они тесно связаны с построениями систем обнаружения вторжений (Network Intrusion Detection). Данные задачи становится сложнее из-за роста различных приложений, распространения зашифрованных протоколов и общего роста трафика. За счет чего возрастает сложность традиционных алгоритмов, например таких как классификация на основе блоков данных(Payload-Based Classification).


%(Кратко про Payload-Based Classification или же DPI ((Deep Packet Inspection ) Базируется на содержании трафика, поэтому не подходит для зашифрованного, также есть модификации, но подходят только для HTTPS трафика.)

Задача управления интернет-трафиком является актуальной в условиях широкого распространением интернета среди населения всего мира. Особенно эта проблема обостряется начиная с 2010-х годов в связи с ростом мирового объема трафика. На это указывают многие научные научные работы, например согласно данным одного из них за последние 10 лет объем трафика возрос почти в 10 раз \cite{reva}. Также на это указывают и сами телекоммуникационные компании. Так например, согласно данным компании <<De CIX>> за период в 5 лет во Франкфурте объем трафика вырос более чем в 2 раза (Рисунок \ref{fig:traffic_growth}) \cite{decix}


\begin{figure}[ht]
	\centerfloat{
		\includegraphics[width=\textwidth]{traffic_growth}
	}
	\caption{Средний и пиковый объемы трафика во Франкфурте в период с 2018 по 2022 год }\label{fig:traffic_growth}
\end{figure}

Общемировым трендом является также рост использования протоколов шифрования и методов скрытия (обфускации) данных и распространение использующих их различных мобильных и десктопных приложений. Так, например, согласно данным компании Google \cite{transparencyreport} процент использования протокола HTTPS на ОС Windows в России на момент конца 2022 года составляет более 93\% (Рисунок \ref{fig:crypto_growth}), а с платформ на ОС Android более 96\%.

% еще про рост шифрования - Sherry, Justine, et al. "Blindbox: Deep packet inspection over encrypted traffic." Proceedings of the 2015 ACM Conference on Special Interest Group on Data Communication. 2015

\begin{figure}[ht]
	\centerfloat{
		\includegraphics[width=\textwidth]{crypto_growth}
	}
	\caption{Рост процента использования HTTPS протокола в трафике, проходящем через сервисы компании Google}\label{fig:crypto_growth}
\end{figure}

Эти факторы приводят к уменьшению эффективности традиционных методов, таких как идентификация на основе портов(port-based identification), классификация на основе блоков данных(payload-based identification) и анализ пакетов на основе их содержания (DPI - deep packet inspection). На эту проблему указывают различные современные исследования \cite{dpi1}, \cite{dpi2}, \cite{dpi3}. В связи с этим возникает необходимость в разработке альтернативных методов и инструменты, основанные на машинном обучении и статистическом анализе являются перспективными и развивающимися \cite{Getman2020}

Еще одним фактором, усложняющим анализ современных сетей является рост трафика IoT сегмента(он же <<интернет вещей>> - использование различных <<умных>> бытовых устройств, обменивающихся данными по сети, Рисунок \ref{fig:IoT}). Рост IoT является быстрым и революционным, но он также представляет значительные проблемы, особенно в области сетевого анализа и безопасности. Сложность и неоднородность сетей IoT, наряду с их быстрым расширением, привели к появлению различных проблем, которые необходимо решать для обеспечения надежной и безопасной работы. IoT расширяется по всему миру, интегрируясь в различные приложения, такие как умные города, сети и военные системы, что увеличивает сложность и неоднородность сети \cite{Chaabouni2019Network, Mo2020Reliability, Wijethilaka2021Survey}. 
Некоторые исследования подчеркивают, что распространение устройств Интернета вещей привело к проблемам в хранении данных, конфиденциальности, протоколах связи и управлении устройствами из-за их ограниченной емкости и мощности \cite{Achkouty2024RDSC}. Также сети IoT уязвимы для кибератак. Примером таковых является инцидент с вредоносным ПО Mirai, который использовал устройства IoT для распределенных атак типа «отказ в обслуживании»(denial-of-service attack, DoS-атака). Это подчеркивает необходимость надежных систем обнаружения сетевых вторжений (NIDS) \cite{Chaabouni2019Network, Patidar2024Cyberattack}. 

\begin{figure}[ht]
	\centerfloat{
		\includegraphics[width=0.8\textwidth]{IoT}
	}
	\caption{Схематичное изображение концепции IoT \cite{iot_pic}}\label{fig:IoT}
\end{figure}

Традиционные методы сетевой мониторинга является неэффективным для IoT из-за разнообразия устройств. Поэтому в настоящее время разрабатываются различные методы, основанные на машинном обучении для распознавания типов устройств посредством анализа сетевого трафика, что позволяет достичь высокой точности в идентификации устройств IoT \cite{Shahid2018IoT}. И поскольку IoT продолжает развиваться и расширяться, разработка эффективных и безопасных решений для управления сетями имеет решающее значение.


Также важность поиска альтернативных моделей оценки сети связана также и с экологическим аспектом. Использование глубоких нейронных сетей, в частности NLP моделей, приводит к значительным выбросам CO2 при их обучении. Согласно исследованиям, использование больших вычислительных мощностей в ходе процесса обучения приводят к большим выбросам CO2 и для некоторых моделей объем выбросов сопоставим с несколькими жизненными циклами автомобилей \cite{Strubell2020-ng}.


%Стоит подчеркнуть, что данная задача важна для различных организации и компании с точки зрения безопасности данных, так как анализ трафика может помочь отслеживать утечки. Пример подобного ...


%Мы можем сделать \textbf{жирный текст} и \textit{курсив}.

\section{Обзор существующих решений анализа интернет трафика}\label{sec:ch1/sec2}

% кратко про приоритезацию трафика

В последние несколько лет фокус исследований науки о данных сместился в сторону использования глубоких нейронных сетей (DNN), и это можно наблюдать во многих областях \cite{TIAN2020251, Rezaei2019, Ferdowsi2019}. Эту тенденцию можно легко объяснить: главное преимущество DNN заключается в том, что они могут выполнять некоторое неявное извлечение признаков самостоятельно, без необходимости в глубоком анализе специализированных инженеров по извлечению признаков из данных. Однако стоит подчеркнуть, что для создания эффективной нейронной сети, инженер должен разработать ее архитектуру, а это уже довольно трудоемкая задача. Главная же проблема заключается в том, что DNN имеют более высокую вычислительную сложность по сравнению с другими алгоритмами, в частности которые включают явное извлечение признаков. Однако стоит подчеркнуть, что последние достижения в разработке оборудования могут позволить смягчить недостатки DNN в вопросе высокой вычислительной сложности. Эти достижения позволят подчеркнуть главный плюс DNN, который заключается в том, что они обеспечивают высокую точность для широкого спектра задач \cite{Liu2019}. 


Одновременно с этим классические методы анализа сетевого трафика, включая такие подходы, как глубокий анализ пакетов (Deep Packet Inspection, DPI), постепенно утрачивают свою эффективность в связи с ростом объёмов трафика, массовым распространением криптографических протоколов и динамичной эволюцией сетевой инфраструктуры, включающей появление новых протоколов передачи данных. Как отмечается в работе \cite{Sherry2015}, увеличение доли зашифрованного трафика делает невозможным применение DPI для детального анализа содержимого пакетов, что существенно снижает применимость данного метода в современных сетях.

Тем не менее, на современном этапе развития информационных технологий возникает острая потребность в разработке и внедрении методов обработки данных, в частности интернет-трафика и данных, генерируемых вредоносным программным обеспечением, с учётом жёстких аппаратных ограничений. Особый интерес представляют сценарии, в которых обработка осуществляется на устройствах с ограниченными вычислительными ресурсами, такими как промежуточные сетевые маршрутизаторы или даже роутеры. В таких условиях традиционные подходы, основанные на глубоких нейронных сетях (DNN), зачастую оказываются неприменимыми из-за высокой вычислительной сложности и значительных требований к памяти и энергопотреблению.

В этом контексте значительное внимание привлекают алгоритмы «классического» машинного обучения, которые, хотя и могут уступать DNN по точности, обладают существенно более низкой вычислительной сложностью и меньшими требованиями к ресурсам. Такие свойства делают их особенно перспективными для использования в системах реального времени, где критически важны быстродействие и энергоэффективность — например, в задачах прогнозирования аварий на транспортных магистралях \cite{Athanasios2019} или активного шумоподавления в аудиосистемах \cite{Xiao2018}. Спектр приложений, ограниченных вычислительными, временными и энергетическими ресурсами, является широким и продолжает активно расширяться, что подчёркивает значимость таких задач по сравнению с теми, что решаются с использованием глубокого обучения.

%%%%%%%%%%%%%%%%%%%%%%%%%%%%

Рассмотрим подробнее недостатки старых методов. Классические методы анализа сети, такие как DPI, port-based-analysis и payload analysis долгое время являлись фундаментальными при управлении сетями и обеспечении их безопасности. Однако в наши дни применение этих становиться все более затруднительным в условия сложности сегодняшних современных сетей.
\cite{Hou2020, Lin2014}

\textbf{Deep Packet Inspection (DPI)}: DPI - это метод, который изучает часть данных (payload, полезную нагрузку) пакета и его заголовок \cite{Chaudhary2011, Oklilas2017}. Данные метод выходит за рамки простой информации о заголовках, позволяя идентифицировать приложения, обнаруживать вторжения в сеть и обеспечивать требуемый уровень качества обслуживания (QoS)\cite{Shubbar2017, Lin2014}.

DPI используется для сопоставления входящего трафика с базой данных сигнатур, представляющих атаки, уязвимости, вирусы и типы трафика.
\cite{Chaudhary2011}


Например, он может идентифицировать конкретные приложения или вредоносные программы, отыскивая уникальные паттерны в полезной нагрузке (payload) пакетов \cite{Lin2014}. DPI является ключевой технологией в системах обнаружения и предотвращения вторжений \cite{Shubbar2017, Parvat2015}


\textbf{Port-Based Analysis}: Этот метод классифицирует сетевой трафик на основе номеров портов, которые традиционно ассоциируются с определенными приложениями или сервисами. Например, порт 80 обычно используется для HTTP-трафика, а порт 21 - для FTP \cite{Hou2020}.

\textbf{Payload Analysis}: Он включает в себя проверку фактического содержимого данных (полезной нагрузки) сетевых пакетов для идентификации приложения или обнаружения вредоносного содержимого. Анализ полезной нагрузки является более глубоким, чем анализ на основе портов, но при этом требует больших вычислительных затрат \cite{Khandait2020}. В отличие от DPI анализирует меньший перечень данных, за счет чего может потреблять меньше вычислительных ресурсов.

Рассмотрим причины, в силу которых данные методы начинают терять свою эффективность в современных сетях:

1.	Развитие и эволюция современных сетей. Современные сетевые сервисы все чаще используют нестандартные порты и обфускацию(то есть маскировку) протоколов, что снижает точность классификации на основе портов \cite{Hou2020}.

2.	Шифрование трафика. Широкое распространение технологий шифрования, таких как TLS (Transport Layer Security), затрудняет проверку полезной нагрузки пакетов (payloads). Традиционные методы DPI основаны на инспекции открытого (незашифрованного) содержимого, что невозможно в случае зашифрованного трафика. Содержимое пакетов скрыто, что делает методы, основанные на полезной(payload) нагрузке, менее эффективными \cite{Ji2024, Shen2023}.


3.	Полиморфные атаки: Полиморфные вредоносные программы и атаки могут менять свои сигнатуры, что затрудняет их обнаружение с помощью DPI на основе заранее определенных шаблонов. Традиционные методы, основанные на сигнатурах, с трудом справляются с этими эволюционирующими угрозами (Рисунок \ref{fig:poly})\cite{Xue2024, Stein2024}


\begin{figure}[ht]
	\centerfloat{
		\includegraphics[scale=0.8]{intro/Polymorphic-virus-structure}
	}
	\caption{Общая схема действия полиморфных атак}\label{fig:poly}
\end{figure}

4. Вычислительная сложность: DPI - это метод, требующий больших вычислительных затрат, поскольку он включает в себя глубокую проверку полезной нагрузки пакетов \cite{Khandait2020, Shubbar2017}. Это может привести к снижению производительности, особенно в высокоскоростных сетях \cite{Shubbar2017}. Сопоставление регулярных выражений (распространенная техника в DPI) может оказаться довольно ресурсоемкой \cite{Nam2022}.

5. Сложность данных: Современные сети генерируют огромные объемы данных, что затрудняет анализ трафика в режиме реального времени с помощью традиционных методов. Данные методы анализа сетей часто не справляются с масштабами и динамикой современного сетевого трафика \cite{Jiang2021}.


6. Проблемы конфиденциальности: Глубокая инспекция пакетов (DPI) предполагает изучение содержимого сетевых соединений, что вызывает серьезные опасения по поводу конфиденциальности \cite{Barut2023}. Пользователи становятся все более чувствительными к последствиям мониторинга сети для конфиденциальности, что усложняет развертывание решений на основе DPI \cite{Belenguer2023}.

7. Атаки на основе искусственного интеллекта. Злоумышленники используют искусственный интеллект для обхода традиционных мер безопасности. Защита систем от таких атак становится все более актуальной \cite{Roshan2024}.


На Рисунке \ref{fig:dpi1} отображена схема этапов обработки сетевых данных, демонстрирующая, как сетевые пакеты проходят через различные %функциональные
блоки, включая парсер, проверку контрольной суммы, менеджер трафика и депарсер, что подчеркивает сложность классического сетевого анализа. В связи с этими проблемами исследователи изучают новые методы, такие как машинное обучение и искусственный интеллект, для повышения эффективности анализа сетевого трафика и безопасности сетей \cite{Shen2023, Ji2024}. Например, машинное обучение можно использовать для выявления закономерностей в зашифрованном трафике без необходимости расшифровки полезной нагрузки \cite{Shen2023}.


\begin{figure}[ht]
	\centerfloat{
		\includegraphics[scale=0.9]{intro/DPI1}
	}
	\caption{Блок-схема этапов работы метода Deep Packet Inspection (DPI)}\label{fig:dpi1}
\end{figure}



\section{Обзор существующих решений оценки производительности сетей}\label{sec:ch1/sec3}

% проблематика - проблемы сетей, ботленки,
Анализ производительности современных сетей сталкивается с рядом проблем, в частности с задачей выявления и устранении различных проблемных узлов, например так называемых <<бутылочных горлышек>> (элемент, ограничивающий пропускную способность пути) (Рисунок \ref{fig:bottleneck}). Традиционные методы часто не оправдывают ожиданий из-за сложности и динамичности современных сетей. Машинное обучение предлагает многообещающие альтернативы для преодоления этих ограничений.
 
 \begin{figure}[ht]
 	\centerfloat{
 		\includegraphics[width=0.8\textwidth]{bottleneck}
 	}
 	\caption{Пример проблемы <<бутылочного горлышка>> в устройстве сетей \cite{bottleneck}}\label{fig:bottleneck}
 \end{figure}

Традиционные методы анализа производительности сети, такие как основанные на теории очередей, часто сталкиваются со сложностью и нестационарными рабочими нагрузками современных сетей, таких как платформы IoT и распределенные системы. Эти методы обычно требуют обширных человеческих знаний и подходов проб и ошибок, которые могут быть неэффективными и подверженными ошибкам \cite{Ouedraogo2022Adaptive, Alsayasneh2020On}. Кроме того, они не могут эффективно обрабатывать стохастическую природу сетевых весов или динамические взаимодействия в многоуровневых распределенных системах \cite{Alsayasneh2020On, Åkerblom2021Online}.


Методы машинного обучения предоставляют инновационные решения для анализа производительности сети и выявления узких мест. Например, метод опорных векторов (SVM) использовался для обнаружения аномалий производительности и выявления узких мест в облачных центрах обработки данных, предлагая высокую точность и низкие накладные расходы по сравнению с традиционными статистическими методами \cite{El_Shamy2021Anomaly}. Аналогичным образом, модели машинного обучения, такие как кластеризация k-средних и искусственные нейронные сети, применялись для прогнозирования пропускной способности сети и использования ресурсов, демонстрируя высокую точность в выявлении узких мест в высокоскоростных сетевых средах \cite{Chauhan2017Applying}.


%Несколько исследований были сосредоточены на выявлении и управлении узкими местами сети с помощью машинного обучения. Например, было предложено использовать сэмплирование Томпсона (Thompson Sampling) в рамках комбинаторной задачи полубандита (semi-bandit problem) для выявления узких мест в сетях со стохастическими весами, что обеспечивает приблизительное решение вычислительно неразрешимой проблемы (невозможно найти решение за приемлемое количество времени или ресурсов) \cite{Åkerblom2021Online}. Кроме того, разрабатываются и  методы адаптивного анализа производительности для оптимизации выявления узких мест в платформах Интернета вещей (IoT), минимизируя дополнительные расходы на мониторинг и сохраняя высокую производительность анализа \cite{Ouedraogo2022Adaptive}.


Несмотря на достигнутые результаты, остаются проблемы в эффективном внедрении методов машинного обучения в системы анализа производительности сети. Сложность сетевых сред, необходимость сбора данных в реальном времени и интеграция с существующей инфраструктурой являются значительными препятствиями \cite{Ouedraogo2022Adaptive, Mai2015Optimizing}. Поэтому дальнейшую работу следует сосредоточить на повышении эффективности методов в сравнении с нейросетевыми методами, в частнсоти  графовыми нейронными сетями  (sampling techniques in graph neural networks), и улучшении совместимости решений машинного обучения с текущими аппаратными и программными системами \cite{Mai2015Optimizing, Wang2021Empirical}.

%\textbf{Сравнение методов}
%\textbf{Network Performance Analysis: Latency, Throughput, and Bottleneck Detection}

Рассмотрим более детально различные подходы в анализе сетевой инфраструктуры. Время отклика (или же задержка), пропускная способность (скорость передачи данных) и наличие в сети проблемных узлов являются фундаментальными показателями производительности в компьютерных сетях. Исторически они отслеживались с помощью систем, основанных на правилах (rule-based systems), и таких протоколов, как SNMP (Simple Network Management Protocol), а также с помощью выборки пакетов (NetFlow/sFlow/IPFIX) и активных измерений (чаще всего - ping и traceroute). SNMP опрашивает счетчики устройств (процессор, байты интерфейса и т. д.) и использует пороговые значения для выявления проблем, а сборщики потоков (flow samplers) собирают сводки о потоках трафика. Однако эти традиционные подходы становятся все более недостаточными в современных высокоскоростных и динамически меняющихся сетях. Например, SNMP фокусируется на состоянии устройств и «не обеспечивает погружение в паттерны сетевого трафика», а его модель опроса может сама по себе создавать большую нагрузку в больших сетях \cite{snmp1}. Аналогично, статические пороговые значения, основанные на правилах, не могут адаптироваться к изменяющемуся трафику: они часто приводят к перегрузке операторов и реагируют только после возникновения проблем \cite{mlmon1}. В сложных современных сетях (с виртуализацией, облачными сервисами, IoT 4G/5G и т. д.) простые опросы или пороговые оповещения могут не обнаружить тонких перегрузок или аномалий. Как отмечается в одном из аналитических материалов, «традиционный мониторинг сети в значительной степени зависит от вмешательства человека и систем, основанных на правилах (rule-based systems)», что приводит к медленному времени реагирования и трудностям с выявлением аномалий в динамическом трафике \cite{mlmon1}. Поэтому администраторы сетей  чаще обращаются к более богатой телеметрии (такие как экспорт потоков, захват пакетов) и аналитике.

%\textbf{Traditional vs. Modern Monitoring: Why Old Methods Falter}

Современные сети генерируют огромные потоки разнородных данных (например, от устройств IoT и пользователей сотовой связи). В таких условиях традиционные методы не справляются с масштабом и сложностью \cite{Abbasi2021}. Например, опрос миллиона датчиков IoT по протоколу SNMP перегрузит как сами устройства, так и сервер управления, а фиксированные пороговые значения не смогут уловить нестационарные характеристики трафика. В исследованиях отмечается, что «традиционные методы управления сетью для мониторинга и анализа данных сталкиваются с проблемами в эффективной обработке больших данных в режиме реального времени», особенно когда структура трафика усложняется за счет мобильности и неоднородности. Мониторинг на основе потоков (NetFlow/sFlow) частично смягчает эту проблему, агрегируя трафик вместо постоянных опросов; он обеспечивает «более глубокую видимость моделей сетевого трафика» и лучше масштабируется \cite{snmp1}. Тем не менее, даже данные о потоках собираются выборочно и могут упускать короткие всплески или малозаметные проблемы с производительностью. По сути, статическая аналитика, основанная на правилах (с использованием счетчиков SNMP или фиксированных формул), теряет свою актуальность в связи с объемом и динамичностью современного трафика.

Для обнаружения узких мест в сети традиционно используются такие метрики, как загрузка канала или длина очереди: например, определение «максимально загруженного» канала по счетчикам интерфейса SNMP или поиск узла с наибольшей задержкой. Однако в высокоскоростных сетях или распределенных системах узкие места могут быстро смещаться и не быть заметными на каком-либо отдельном устройстве. Существуют такие методы, как сетевая томография (определение задержек в каналах на основе сквозных зондов) и профилирование на основе трассировки, но они сложны и часто работают в автономном режиме. На практике инженеры все чаще дополняют SNMP сложной аналитикой или специализированной телеметрией (например, внутриполосной телеметрией, данными P4). Несмотря на эти достижения, системы мониторинга все еще часто полагаются на фиксированные эвристические правила (например, флаг, срабатывающий при использовании канала >90 \%), которые могут пропустить коррелирующие аномалии между каналами или плохо адаптироваться к изменяющимся нагрузкам.

%\textbf{Monitoring Low-Power/Resource-Constrained Devices}
Развитие IoT и пограничных вычислений создает новые проблемы. Сети IoT состоят из «большого количества разнородных и ограниченных в ресурсах устройств (с точки зрения памяти, процессора и пропускной способности)», работающих на беспроводных каналах с потерями \cite{Aboubakar2022}. Эти устройства часто находятся в состоянии сна для экономии энергии, имеют ограниченную память и пропускную способность сети. В результате традиционные протоколы мониторинга оказываются слишком тяжеловесными: агенты SNMP и частые опросы могут разрядить батарею или превысить пропускную способность устройства. Для решения этой проблемы были разработаны облегченные протоколы управления. Например, протокол OMA Lightweight M2M (LwM2M) и его адаптации, такие как 6LoWPAN-SNMP, направлены на расширение SNMP-подобного мониторинга для IPv6-сетей с низким энергопотреблением \cite{Aboubakar2022}. Эти протоколы сжимают или упрощают SNMP-сообщения и работают по UDP/CoAP вместо TCP. Другие усилия включают управление на основе CoAP (например, CoMI) и модели «тяни-толкай», которые позволяют устройствам отправлять критические обновления только по определенным событиям.

Но даже в этом случае мониторинг должен быть тщательно продуман. В одном из исследований подчеркивается, что устройства IoT «крайне ограничены по мощности, памяти и пропускной способности», поэтому APM (мониторинг производительности приложений) для IoT должен производиться с минимальными накладные расходами \cite{GRAF2024989}. Это означает, что необходимо проводить сбор только основных показателей, использовать сжатие данных и, возможно, внутрисетевую агрегацию. Например, вместо постоянного экспорта потоков пограничные узлы могут делать локальные выборки или вычислять локальную статистику (среднюю задержку, потерю пакетов и т.д.) и передавать только сводные данные. Также наблюдается тенденция к пограничным вычислениям или обработке на месте: анализ данных на устройстве или на ближайшем шлюзе, так что в облако нужно отправлять только аномалии или сжатые результаты. Таким образом, для мониторинга IoT и крайних узлов требуются технологии с ультранизким энергопотреблением: протоколы и аналитика, которые учитывают ограничения процессора и батареи, но при этом обеспечивают понимание уровня задержки и пропускной способности.

%\textbf{Machine Learning Approaches}
Учитывая описанные проблемы, многие исследователи рассматривают машинное обучение (ML) в качестве нового инструмента для анализа производительности сети. В отличие от фиксированных правил, модели ML могут изучать сложные взаимосвязи на основе данных (трафик, метрики устройств) и адаптироваться с течением времени. Классические методы ML (например, SVM, деревья решений, случайные леса) применялись для решения таких задач, как классификация трафика, обнаружение аномалий и прогнозирование производительности. Например, в одной из работ авторы продемонстрировали, что система на основе SVM может решать несколько задач сетевого анализа более точно и с меньшими затратами на измерения, чем традиционные эвристики. В каждом случае они обнаружили, что решения на основе SVM значительно превосходят существующие решения по точности, затратам на измерения и простоте практического развертывания \cite{Mirza2012}. Такие методы обычно требуют разработки характеристик (например, выбора релевантных счетчиков или статистики потоков), но достаточно легки в вычислениях, чтобы работать на серверах или даже интеллектуальных пограничных шлюзах.

Глубокое обучение (DL - Deep Learning) - это область с применением нейронных сетей с большим количеством слоев. Данный подход также набирает обороты для мониторинга сети (иногда называемого NTMA: мониторинг и анализ сетевого трафика). В обзорах отмечается, что современные системы производят массивные, гетерогенные данные, с которыми традиционная аналитика не справляется \cite{Abbasi2021}. Глубокие сети (такие как CNN, RNN, LSTM, автоэнкодеры) могут автоматически извлекать особенности из данных необработанных пакетов или временных рядов и обнаруживать тонкие закономерности. Исследователи использовали глубокие модели для классификации трафика, прогнозирования объема и обнаружения аномалий. Например, отмечается, что «глубокое обучение эффективно используется для облегчения аналитики и обнаружения аномалий в системах больших данных для распознавания скрытых и сложных закономерностей», что мотивирует его применение для классификации и прогнозирования трафика \cite{Abbasi2021}. В целом стоит отметить, что на практике ML/AI подход может обеспечить предиктивное оповещение (т.е. прогнозирование перегрузки канала до ее возникновения) и анализ первопричины (сопоставление разрозненных показателей для выявления узкого места). Многие отраслевые решения уже включают в себя конвейеры ML для обнаружения аномалий и оценки пропускной способности.


%\textbf{Limitations of ML and Deep Learning in Low-Power Settings}
Несмотря на свои перспективы, ML и особенно глубокое обучение имеют существенные недостатки в ограниченных условиях. Глубокие нейронные сети (например, CNN) обычно включают от тысяч до миллионов параметров (весов) и требуют значительных ресурсов по CPU, GPU и памяти. Исследования показывают, что, хотя CNN могут достигать высокой точности на сложных данных, но «они требуют огромных вычислительных ресурсов и большого объема данных для эффективного обучения модели», что является очевидным ограничением для IoT-устройств \cite{Jebur2025}. Даже запуск обученной глубокой модели (инференс) может потреблять нетривиальную энергию и задержку, что делает ее непригодной для маломощных датчиков. Классические модели ML (например, деревья решений или небольшие нейронные сети) более легкие, но все равно нуждаются в периодическом переобучении на основе новых данных. В целом, подходы ML меняют подход с использования ресурсов на интеллектуальный анализ: они уменьшают количество ложных тревог, обучаясь закономерностям, но при этом требуют данных для обучения, вычислительных мощностей и памяти \cite{mon_chall}.


Следует отметить, что при недостаточно продуманной архитектуре глубокие модели могут демонстрировать склонность к переобучению.
%или игнорировать многомерные компромиссы (задержка против пропускной способности против использования ресурсов), если они тщательно не разработаны. 
Также в одном из аналитических материалов отмечается, что «компромиссы между улучшением обнаружения аномалий и управлением вычислительными затратами редко обсуждаются», что означает, что модели часто хорошо работают на идеальных лабораторных данных, но не работают в реальных сетях \cite{Jebur2025}. Наконец, как классические, так и глубокие методы ML требуют маркированных данных или четко определенных базовых показателей, которые бывает сложно получить в новых развертываниях. В итоге, несмотря на то, что ML обеспечивает адаптивность и более глубокое понимание, размер модели, вычислительные затраты и требования для работы в реальном времени делают его сложным для развертывания на маломощных или критичных к задержкам устройствах \cite{Jebur2025}.


Каждый подход имеет свои плюсы и минусы. Традиционный мониторинг легок и прост в развертывании, но может лишь реактивно (то есть только через некоторое время после возникновения) отмечать заранее известные типы проблем. Классический ML добавляет «интеллектуальность» без лишних накладных расходов, повышая точность обнаружения и снижая количество ложных срабатываний \cite{Mirza2012}. Глубокое обучение может моделировать чрезвычайно сложное поведение и корреляции, но его требования к ресурсам делают его непрактичным для многих сценариев с применением IoT и краевых устройств (edge devices) \cite{Jebur2025}. На практике часто лучше всего работают гибридные стратегии - например, использование легкого ML на границе для немедленного обнаружения аномалий и отправка обобщенных данных в центр обработки данных для более глубокого анализа. Широкий спектр исследований и обзоров (например, обзоры по ML для сетевых технологий \cite{Abbasi2021, Mirza2012}, управлению сетями для IoT \cite{Aboubakar2022, GRAF2024989} и современным тенденциям мониторинга \cite{snmp1, mlmon1}) подтверждает, что классический протокол SNMP и мониторинг на основе правил испытывают трудности в высокопроизводительных, динамичных сетях, и подчеркивает необходимость в адаптивных, эффективных методах. В то же время публикации по ML для мониторинга сетей \cite{Mirza2012, Jebur2025} подчеркивает как перспективы, так и ограничения подходов, основанных на искусственном интеллекте. Краткие итоги сравнительного анализа приведены в Таблице \ref{tab:methods1}. В итоге стоит подчеркнуть, что ни один из данных подходов не является оптимальным во всех аспектах. В частности на сегодняшний распространено применение DL для различных областей сетевого анализа, однако будущее сетевого мониторинга видится в применении легковесных и  оптимизированных как DL так и ML алгоритмов, позволяющие их использование на маломощных устройствах.

Итак, можно резюмировать то, что в то время как традиционные методы анализа  сети имеют различные ограничения, машинное обучение предлагает многообещающие альтернативы для выявления и управления проблемными сегментами. Проведение исследований и разработок в этой области имеет важное значение для решения вышеописанных задач и полного использования потенциала методов машинного обучения в современных сетевых средах.

\begin{comment}
\begin{table}[htbp]
	\centering
	\caption{Сравнение подходов к способам мониторинга сетей}
	\label{tab:methods1}
	\begin{tabular}{|p{3cm}|p{3cm}|p{3cm}|p{3cm}|p{3cm}|}
		\hline
		\textbf{Подход}  & \textbf{Гибкость} & \textbf{Вычислительная сложность} & \textbf{Точность в комплексных сценариях} & \textbf{IoT и граничные вычисления}  \\ \hline
		Традиционный (SNMP, пороговые правила, выборка потоков) & Статический (фиксированные пороги, ручные правила). & Низкая выч. нагрузка (простые опросы/ пороговые правила). & Приемлемая для известных проблем, но плохая  для новых или сложных аномалий. & Хорошо подходит для высокоуровневого мониторинга (серверы, ЦОД), плохо для IoT (из-за стоимости опроса). \\ \hline
		Классический ML(SVM, decision tree, k-NN, и т.д.) & Умеренная (модели могут адаптироваться при переобучении). & Средняя (для обучения требуется ЦП; вывод (инференс) малопотребляющий). & Хорошая во многих случаях - значительно лучше, чем статические правила \cite{Mirza2012}.& Умеренно эффективно (может выполняться в облаке или на пограничных шлюзах). \\ \hline 
		Глубокое обучение(CNN, RNN, DNN и т.д.) & Высокая (способен постоянно адаптироваться к новым данным и использовать обучение без учителя). & Высокая (крупные модели, рекомендуется использование GPU/TPU; высокая нагрузка на оборудование). & Отлично работает с большими и сложными наборами данных (обнаруживает тонкие аномалии и нелинейные зависимости).\cite{Abbasi2021} & Плохо подходит для маломощных устройств . Технологии TinyML могут помочь.  \\ \hline
	\end{tabular}
\end{table}
\end{comment}

\begin{table}[htbp]
	\centering
	\caption{Сравнение подходов к способам мониторинга сетей}
	\label{tab:methods1}
	\begin{tabular}{|>{\centering\arraybackslash}p{3.5cm}|>{\centering\arraybackslash}p{4.6cm}|>{\centering\arraybackslash}p{3.5cm}|>{\centering\arraybackslash}p{4cm}|}
		\hline
		\textbf{Подход}  & \textbf{Вычислительная сложность} & \textbf{Точность в комплексных сценариях} & \textbf{IoT и граничные вычисления}  \\ \hline
		Традиционный (SNMP, пороговые правила, выборка потоков)  & Низкая нагрузка (простые опросы/ пороговые правила). & Приемлемая для известных проблем, но плохая  для новых или сложных аномалий. & Хорошо подходит для высокоуровневого мониторинга (серверы, ЦОД), плохо для IoT (из-за стоимости опроса). \\ \hline
		Классический ML(SVM, decision tree, k-NN, и т.д.) & Средняя (для обучения требуется ЦП; инференс малопотребляющий) & Хорошая во многих случаях. Значительно лучше, чем статические правила \cite{Mirza2012}.& Умеренно эффективно (может выполняться в облаке или на пограничных шлюзах). \\ \hline 
		Глубокое обучение(CNN, RNN, DNN и т.д.) & Высокая (крупные модели, рекомендуется использование GPU/TPU; высокая нагрузка на оборудование). & Отлично работает с большими и сложными наборами данных \cite{Abbasi2021} (обнаруживает тонкие аномалии и нелинейные зависимости). & Плохо подходит для маломощных устройств Технологии TinyML могут помочь.  \\ \hline
	\end{tabular}
	
\end{table}



\FloatBarrier
